\documentclass[10pt,a4paper]{amsart}
\usepackage[margin=19mm]{geometry}
\usepackage{amsmath,amssymb,mathtools}
\usepackage[scr=rsfs,cal=euler]{mathalfa}
\usepackage{xfrac}
\usepackage[unicode,hidelinks,bookmarks=false]{hyperref}
\newcommand{\ie}{i.e.\ }
\newcommand{\cf}{cf.\ }
\newcommand{\resp}{respectively\ }
\newcommand{\Subsection}{\subsection}
\providecommand{\nobreakdash}{\nobreak\hbox{-}}
\setlength{\emergencystretch}{2em}
\multlinegap0pt
\usepackage{bm}
\usepackage{bbm}
\usepackage{dsfont}
\usepackage{xspace}
\usepackage{cancel}
\usepackage{mathtools}
\usepackage{amsthm}
\makeatletter
\@ifundefined{prop}{\newtheorem{prop}{Prop}}{}
\@ifundefined{lem}{\newtheorem{lem}[prop]{Lemma}}{}
\makeatother
\newcommand{\Sph}{\mathbb{S}}
\newcommand{\Sym}{\mathfrak{S}}
\title[Dirichlet-Voronoi domain and injectivity radius of flag mani]{Dirichlet-Voronoi domain and injectivity radius of flag manifolds -- equivariant cell structure on $O(3)/O(1)^3$}
\author{Arthur Garnier}
\date{Source: arXiv:2011.06338v5 (CC BY 4.0)}
\begin{document}
\maketitle

\noindent\emph{Source and licence.} Dirichlet-Voronoi domain and injectivity radius of flag manifolds -- equivariant cell structure on $O(3)/O(1)^3$. Version arXiv:2011.06338v5 (CC BY 4.0), distributed under the Creative Commons Attribution 4.0 International licence, as recorded in the arXiv licence metadata for this exact version and shown on the abstract page https://arxiv.org/abs/2011.06338v5. Full rights notice: licenses/arxiv-2011.06338-CC-BY-4.0.txt.\par\smallskip

\noindent\textit{Editorial scope.} Counted unit: the Lem whose source label is \texttt{tcube} in the article (source0.tex lines 1535--1547, source label \texttt{tcube}), delivered together with the article's own complete original proof of it (source0.tex lines 1548--1582, one \texttt{proof} environment of 35 lines). No mathematics is written, repaired or re-derived: statement and proof are the article's own text line for line. Required context retained uncounted: none. Every definition, display and label of those lines is retained unchanged. Omitted from this package and not redistributed: 1--1534, 1583--2655. Nothing inside a retained excerpt is omitted or altered: every retained body is delivered exactly as the article's lines are written, so no omission receipt of any kind accompanies this package. Citations: the retained excerpts carry no citation command at all, so no bibliography entry of the article is transcribed and no reference is redistributed. Reference adaptation: none was needed and none was made; every reference inside the delivered unit resolves inside it, so no unresolved reference or citation is produced. Printed slips kept literal: no printed word, symbol, hypothesis or number of the counted statement or of its proof was altered. Layout and class: the article's own class is \texttt{amsart} with packages including babel, csquotes, amsfonts, hyperref, etoolbox, amsmath, amssymb, amsthm, stmaryrd, dsfont, pifont, caption, subcaption, mathrsfs; the wrapper is the lane's pinned \texttt{amsart} editorial wrapper at 10pt on A4, which restates the theorem-like environments the retained text needs and adds the article's own macro definitions that the retained text needs (\texttt{\textbackslash Sph}, \texttt{\textbackslash Sym}), with extra wrapper packages (bm, bbm, dsfont, xspace, cancel, mathtools). The delivered numbering is the unit's own. Assets: no retained span contains a figure environment, a graphics-inclusion command, a PDF-page inclusion, a table or a TikZ picture; no image file is copied into this package, the package declares no asset dependency and no image file is redistributed with it. Licence: version arXiv:2011.06338v5 of this article is distributed under the Creative Commons Attribution 4.0 International licence, as recorded in the arXiv licence metadata for this exact version and shown on the abstract page https://arxiv.org/abs/2011.06338v5. A full rights notice is delivered at licenses/arxiv-2011.06338-CC-BY-4.0.txt. Characters: every retained excerpt is pure ASCII, so the delivered engine, XeLaTeX with its default Latin Modern text font, reports no missing character.
\begin{lem}\label{tcube}
Let $q=a+bi+cj+dk\in\Sph^3$ be such that $a\ge|b|,|c|,|d|$. Then
\[\Pi(q)\in\mathcal{DV}_3~\Longleftrightarrow~\left\{\begin{array}{rcl}|b|,|c|,|d| & \le & a(\sqrt{2}-1), \\[.5em] |b\pm c\pm d| & \le & a.\end{array}\right.\]
Moreover, letting $s_\gamma:=s_{\alpha+\beta}=s_\alpha s_\beta s_\alpha=s_\beta s_\alpha s_\beta$, we have
\begin{equation}\label{partial}\tag{$\partial$}
\left\{\begin{array}{rcl}
\Pi(q)\in Z_{s_\alpha} & \Longleftrightarrow & a(\sqrt{2}-1)=|d|, \\[.5em]
\Pi(q)\in Z_{s_\beta} & \Longleftrightarrow & a(\sqrt{2}-1)=|b|, \\[.5em]
\Pi(q)\in Z_{s_{\alpha+\beta}} & \Longleftrightarrow & a(\sqrt{2}-1)=|c|, \\[.5em]
\Pi(q)\in Z_{s_\alpha s_\beta} & \Longleftrightarrow & a=\max(b-c-d,-b+c-d,-b-c+d,b+c+d), \\[.5em]
\Pi(q)\in Z_{s_\beta s_\alpha} & \Longleftrightarrow & a=\max(b+c-d,b-c+d,-b+c+d,-b-c-d).\end{array}\right.
\end{equation}
\end{lem}

\begin{proof}
The hypothesis $a\ge|b|,|c|,|d|$ ensures that
\[d(1,\Pi(q))=4\sqrt{3}\arccos\|q\|_\infty=4\sqrt{3}\arccos(a).\]
As the elements $s_\alpha$, $s_\beta$, $s_{\alpha+\beta}=s_\alpha s_\beta s_\alpha$, $s_\alpha s_\beta$ and $s_\beta s_\alpha$ of $\Sym_3$ may be represented in $\Sph^3$ respectively by $\tfrac{1}{\sqrt{2}}(1+k)$, $\tfrac{1}{\sqrt{2}}(1+i)$, $\tfrac{1}{\sqrt{2}}(1+j)$, $\tfrac12(1+i+j+k)$ and $\tfrac12(1-i+j+k)$, we find
\[\left\{\begin{array}{rcl}
d(1,\Pi(q)s_\alpha) & = & 4\sqrt{3}\arccos\left(\tfrac{1}{\sqrt{2}}\max(|a\pm d|,|b\pm c|)\right), \\[.5em]
d(1,\Pi(q)s_\beta) & = & 4\sqrt{3}\arccos\left(\tfrac{1}{\sqrt{2}}\max(|a\pm b|,|c\pm d|)\right), \\[.5em]
d(1,\Pi(q)s_{\alpha+\beta}) & = & 4\sqrt{3}\arccos\left(\tfrac{1}{\sqrt{2}}\max(|a\pm c|,|b\pm d|)\right), \\[.5em]
d(1,\Pi(q)s_\alpha s_\beta) & = & 4\sqrt{3}\arccos\left(\tfrac12\max_{\epsilon_i=\pm1,~\epsilon_1\epsilon_2\epsilon_3=1}(|a+\epsilon_1b+\epsilon_2c+\epsilon_3d|)\right), \\[.5em]
d(1,\Pi(q)s_\beta s_\alpha) & = & 4\sqrt{3}\arccos\left(\tfrac12\max_{\epsilon_i=\pm1,~\epsilon_1\epsilon_2\epsilon_3=-1}(|a+\epsilon_1b+\epsilon_2c+\epsilon_3d|)\right).\end{array}\right.\]
%d(1,\Pi(q)s_\alpha s_\beta) & = & 4\sqrt{3}\arccos\left(\tfrac12\max(|a+b-c-d|,|a-b+c-d|,|a-b-c+d|,|a+b+c+d|)\right), \\[.5em]
%d(1,\Pi(q)s_\beta s_\alpha) & = & 4\sqrt{3}\arccos\left(\tfrac12\max(|a+b+c-d|,|a+b-c+d|,|a-b+c+d|,|a-b-c-d|)\right).\end{array}\right.\]
Thus, the conditions $d(1,\Pi(q))\le d(w,\Pi(q))$ for $w\in\Sym_3\setminus\{1\}$ translate into the system of inequalities
\begin{equation}\label{sy}\tag{$S$}
\left\{\begin{array}{rcl}
a\sqrt{2} & \ge & \max(|a\pm d|,|b\pm c|,|a\pm b|,|c\pm d|,|a\pm c|,|b\pm d|), \\[.5em]
2a & \ge & \max(|a\pm b\pm c\pm d|).\end{array}\right.
\end{equation}
The second set of inequalities in \eqref{sy} is readily equivalent to the inequalities $|b\pm c\pm d|\le a$. On the other hand, since $|b|\le a$, we have $a\pm b\ge 0$ and therefore, the first set of inequalities in \eqref{sy} implies in particular that $a\sqrt{2}\ge a\pm b$, so that $a(\sqrt{2}-1)\ge |b|$. In the same way, $a(\sqrt{2}-1)\ge |c|,|d|$. The triangular inequality then yields $|b\pm c|\le|b|+|c|\le 2a(\sqrt{2}-1)<a\sqrt{2}$ and similarly, $|c\pm d|,|b\pm d|<a\sqrt{2}$. Therefore, the system \eqref{sy} is indeed equivalent to the system of the statement.

Now, using the formulae for the distance $d(1,\Pi(q)w)$ given above, we arrive at
\[\Pi(q)\in H_{s_\alpha}~\Longleftrightarrow~a\sqrt{2}=\max(|a\pm d|,|b\pm c|)=a+|d|~\Longleftrightarrow~a(\sqrt{2}-1)=|d|,\]
the conditions for $H_{s_\beta}$ and $H_{s_{\alpha+\beta}}$ being analogous. 
Observe that since $a\pm b\ge0$, we have the inequality $a-b-c+d\ge -a+b-c+d$ and similarly, $a+b-c-d\ge -a-b-c-d$, $a-b+c-d\ge -a+b+c-d$ and $a+b+c+d\ge -a-b+c+d$, so that
\begin{align*}
&\max(|a+b-c-d|,|a-b+c-d|,|a-b-c+d|,|a+b+c+d|) \\
=&\max(a+b-c-d,a-b+c-d,a-b-c+d,a+b+c+d) \\
=&a+\max(b-c-d,-b+c-d,-b-c+d,b+c+d).
\end{align*}
Thus, using the formula for $d(1,\Pi(q)s_\alpha s_\beta)$, we have equivalences
\begin{align*}
\Pi(q)\in H_{s_\alpha s_\beta}&\Longleftrightarrow 2a=\max(|a+b-c-d|,|a-b+c-d|,|a-b-c+d|,|a+b+c+d|) \\ 
&\Longleftrightarrow a=\max(b-c-d,-b+c-d,-b-c+d,b+c+d).
\end{align*}
\end{proof}

\end{document}
